\section{Discussion}
\label{sec:discussion}

\paragraph{Substrate matters.} All our results concern \emph{weight} decomposition. \bsf's original result is on \emph{activation} decomposition. Language-model activations are known to be non-aligned with principal directions of variance in practice (this is why SAEs beat PCA baselines on interpretability metrics), so we do not claim \svdomp\ would beat \bsf\ on activations. What we do claim: on weight decomposition, the analytic SVD dictionary dominates trained baselines; and the stable-rank plateau that \bsf\ interprets as "concepts are 2--4-dimensional" is method-independent, because analytic blocks reproduce it exactly.

\paragraph{Efficiency angle.} A concurrent benefit of \svdomp\ is that it eliminates the training compute of per-matrix probes. For interpretability at scale (hundreds of layers of a reasoning model, thousands of interpreted modules), the difference between one SVD per module (subsecond on CPU) and a full AdamW training loop per module ($\sim$40--200 steps on GPU) becomes materially significant. On our 24-module comparison the analytic methods complete in seconds; the trained methods dominate wall-clock. We view this as a favorable side effect rather than the primary contribution.

\paragraph{Where trained probes should be used.} Two regimes remain where trained probes are appropriate on weight matrices: (i) non-Frobenius objectives that compose $W$ with downstream nonlinearities and matrices, where our scaffold-mode trainable correction (2 parameters per block) captures 73.9\% of the achievable gain relative to analytic; and (ii) the six specific modules with compressed singular spectra where Davis-Kahan predicts \svdomp\ instability. For these cases, the SVD is the correct \emph{starting point}, not a competitor.

\paragraph{Limitations.} (i) All Goodfire-67M experiments use a small pretrained checkpoint. The Pareto dominance and the stable-rank plateau reproduce on real Pythia-70M activations, but a systematic study on modern reasoning-scale models remains future work; the efficiency argument suggests it is precisely at large scale that \svdomp's zero training cost becomes materially valuable. (ii) The non-Frobenius escape (Sec.\ \ref{sec:non-frobenius-exp}) uses adversarial synthetic $W_\text{next}$; on real neighboring-layer matrices the gain is smaller because $W_\text{next}$ has less structured misalignment with $W$'s SVD, though it remains positive on 24 / 24 modules.

\paragraph{Implications for alignment.} If trained interpretability probes on weight matrices are approximating a quantity that the SVD already provides, three implications follow for the alignment research program. First, when interpretability results depend on probe-specific claims (e.g., "the trained probe found this concept"), those claims should be checked against the analytic SVD baseline before being credited to the training procedure. Second, the actual novelty of a trained probe is what it adds on top of the analytic baseline: the correction, not the base decomposition. Third, at reasoning-model scales, interpretability infrastructure that requires no training per layer becomes tractable in ways that trained probes do not.

\paragraph{Conclusion.} On weight decomposition of a widely used interpretability benchmark model, the SVD dictionary provides a Frobenius-optimal per-input sparse decomposition that dominates trained baselines on 24 of 24 per-metric comparisons and 18 of 24 modules Pareto-overall. \bsf's cross-modal 2--4-dimensional concept finding reproduces on real language-model activations without training. Trained decomposition adds value on non-Frobenius composed objectives; our small trainable correction on top of the SVD scaffold captures 73.9\% of that value. Together, these results position analytic decomposition as the appropriate baseline for interpretability of large neural network weights, and provide a principled account of the regime where training legitimately adds signal.
